% theory_coha_e1_sector.tex
% Working note: The critical CoHA as the E_1-sector of the quantum vertex chiral group
% Raeez Lorgat, April 2026
%
% Develops the precise relationship between the critical CoHA (associative/E_1)
% and the full quantum vertex chiral group (braided/E_2), with complete proofs
% for C^3 and the general toric case.

\documentclass[11pt]{article}
\usepackage[T1]{fontenc}
\usepackage[utf8]{inputenc}
\usepackage{amsmath,amssymb,amsthm}
\usepackage{mathtools}
\usepackage{tikz-cd}
\usetikzlibrary{positioning,arrows.meta}
\usepackage[margin=1.25in]{geometry}
\usepackage{enumitem}

\newtheorem{theorem}{Theorem}[section]
\newtheorem{proposition}[theorem]{Proposition}
\newtheorem{lemma}[theorem]{Lemma}
\newtheorem{corollary}[theorem]{Corollary}
\newtheorem{conjecture}[theorem]{Conjecture}
\theoremstyle{definition}
\newtheorem{definition}[theorem]{Definition}
\newtheorem{example}[theorem]{Example}
\newtheorem{construction}[theorem]{Construction}
\theoremstyle{remark}
\newtheorem{remark}[theorem]{Remark}
\newtheorem{warning}[theorem]{Warning}

\newcommand{\ZZ}{\mathbb{Z}}
\newcommand{\CC}{\mathbb{C}}
\newcommand{\QQ}{\mathbb{Q}}
\newcommand{\NN}{\mathbb{N}}
\newcommand{\Spec}{\operatorname{Spec}}
\newcommand{\Crit}{\operatorname{Crit}}
\newcommand{\Tr}{\operatorname{Tr}}
\newcommand{\Rep}{\operatorname{Rep}}
\newcommand{\Hom}{\operatorname{Hom}}
\newcommand{\End}{\operatorname{End}}
\newcommand{\Sym}{\operatorname{Sym}}
\newcommand{\id}{\operatorname{id}}
\newcommand{\HH}{\operatorname{HH}}
\newcommand{\HC}{\operatorname{HC}}
\newcommand{\CoHA}{\mathcal{H}}
\newcommand{\cA}{\mathcal{A}}
\newcommand{\cC}{\mathcal{C}}
\newcommand{\cF}{\mathcal{F}}
\newcommand{\cM}{\mathcal{M}}
\newcommand{\cO}{\mathcal{O}}
\newcommand{\cZ}{\mathcal{Z}}
\newcommand{\cR}{\mathcal{R}}
\newcommand{\cW}{\mathcal{W}}
\newcommand{\frakg}{\mathfrak{g}}
\newcommand{\frakn}{\mathfrak{n}}
\newcommand{\frakh}{\mathfrak{h}}
\newcommand{\frakb}{\mathfrak{b}}
\newcommand{\Ran}{\operatorname{Ran}}
\newcommand{\Fact}{\operatorname{Fact}}
\newcommand{\op}{\operatorname{op}}
\newcommand{\BM}{\operatorname{BM}}
\newcommand{\DT}{\operatorname{DT}}
\newcommand{\GL}{\operatorname{GL}}
\newcommand{\gl}{\mathfrak{gl}}

\title{The Critical CoHA as the $E_1$-Sector\\
of the Quantum Vertex Chiral Group}
\author{Raeez Lorgat}
\date{April 2026 --- Working Note}

\begin{document}
\maketitle

\begin{abstract}
For a Calabi--Yau threefold $X$ with quiver-with-potential $(Q,W)$, the critical
cohomological Hall algebra $\CoHA(Q,W)$ is an associative ($E_1$) algebra built from
Borel--Moore homology of vanishing cycle sheaves on the critical locus.  We develop
the precise statement and proof that $\CoHA(Q,W)$ is isomorphic to the universal
enveloping algebra $U(\frakn_+)$ of the positive part of the Borcherds--Kac--Moody /
Yangian / quantum vertex chiral group $G(X)$.  The $E_1$-structure (associativity)
is the ordered sector; the passage to the full $E_2$-structure (braiding) adds the
quantum group $R$-matrix via the Drinfeld double construction.  We verify this
completely for $\CC^3$ (Schiffmann--Vasserot: $\CoHA = Y^+(\widehat{\gl}_1)$,
Drinfeld double $= Y(\widehat{\gl}_1) \cong \cW_{1+\infty}$) and for general toric
CY3 without compact 4-cycles (Rapcak--Soibelman--Yang--Zhao).
\end{abstract}

\tableofcontents

%% ======================================================================
\section{Overview and motivation}
\label{sec:overview}
%% ======================================================================

The purpose of this note is to make precise the slogan
\begin{center}
\emph{The critical CoHA is the $E_1$-sector (positive half) of the quantum vertex
chiral group.}
\end{center}

\noindent
This slogan has three layers of content:

\begin{enumerate}[label=(\alph*)]
  \item \textbf{Algebraic.}  The CoHA $\CoHA(Q,W)$ is an associative ($E_1$)
    algebra; the quantum vertex chiral group $G(X)$ is an $E_2$-chiral algebra
    whose representation category is braided monoidal.  The CoHA recovers the
    positive half $U(\frakn_+)$.
  \item \textbf{Operadic.}  The passage from $E_1$ to $E_2$ adds a braiding.
    By Dunn additivity, $E_2 \simeq E_1 \otimes E_1$; concretely, the second
    $E_1$-direction is the Cartan/coweight direction of the triangular decomposition,
    and the braiding is the $R$-matrix.
  \item \textbf{Geometric.}  The CoHA multiplication comes from correspondences
    on moduli of quiver representations (genus-0, tree-level curve counting);
    the $E_2$-enhancement incorporates wall-crossing (quantum dilogarithm
    identities) and loop-level corrections.
\end{enumerate}

\noindent
We work equivariantly with respect to the torus $T = (\CC^*)^3$ acting on CY3
geometry, with equivariant parameters $\epsilon_1, \epsilon_2, \epsilon_3$
satisfying the CY condition $\epsilon_1 + \epsilon_2 + \epsilon_3 = 0$.


%% ======================================================================
\section{The critical cohomological Hall algebra}
\label{sec:critical-coha}
%% ======================================================================

\subsection{Quivers with potential from CY3}
\label{subsec:quiver-potential}

Let $X$ be a toric Calabi--Yau threefold.  The toric data (fan $\Sigma$ in
$\ZZ^3$ with all generators coplanar, satisfying the CY condition) determines
a quiver with potential $(Q, W)$ via the brane tiling / dimer model construction.

\begin{definition}[Quiver with potential]
\label{def:quiver-potential}
A \emph{quiver with potential} $(Q, W)$ consists of:
\begin{itemize}
  \item A quiver $Q = (Q_0, Q_1, s, t)$ with vertex set $Q_0$, arrow set $Q_1$,
    and source/target maps $s, t \colon Q_1 \to Q_0$.
  \item A potential $W \in \CC Q / [\CC Q, \CC Q]$, a formal sum of oriented
    cycles in $Q$ modulo cyclic equivalence.
\end{itemize}
The Jacobian algebra is $\operatorname{Jac}(Q, W) = \CC Q / (\partial_a W : a \in Q_1)$.
The CY3 condition is equivalent to $\operatorname{Jac}(Q,W)$ being a
3-Calabi--Yau algebra in the sense of Ginzburg.
\end{definition}

\begin{example}[$\CC^3$: Jordan quiver]
\label{ex:jordan}
The Jordan quiver has one vertex and one loop $x$.  With three loops
$x, y, z$ and potential $W = \Tr(x[y,z]) = \Tr(xyz - xzy)$, this is the
quiver-with-potential for $\CC^3$.  The representation space of dimension
$d$ is $\operatorname{Mat}_{d \times d}(\CC)^3$ with the function
$W_d = \Tr(X[Y,Z])$ and the $\GL_d$-action by simultaneous conjugation.
\end{example}

\begin{example}[Resolved conifold]
\label{ex:conifold}
The Klebanov--Witten quiver has two vertices and four arrows
$a_1, a_2 \colon 0 \to 1$, $b_1, b_2 \colon 1 \to 0$, with potential
$W = \Tr(a_1 b_1 a_2 b_2 - a_1 b_2 a_2 b_1)$.
\end{example}


\subsection{Definition of the critical CoHA}
\label{subsec:critical-coha-def}

\begin{definition}[Critical CoHA]
\label{def:critical-coha}
For a quiver with potential $(Q, W)$, the \emph{critical cohomological Hall
algebra} is the $\ZZ_{\geq 0}^{Q_0}$-graded vector space
\[
  \CoHA(Q, W) = \bigoplus_{\mathbf{d} \in \ZZ_{\geq 0}^{Q_0}}
    H^{\BM}_*\bigl(\Crit(W_{\mathbf{d}})/\GL_{\mathbf{d}},\,
    \varphi_{W_{\mathbf{d}}}\bigr)
\]
where:
\begin{itemize}
  \item $\mathbf{d} = (d_i)_{i \in Q_0}$ is a dimension vector;
  \item $\GL_{\mathbf{d}} = \prod_{i \in Q_0} \GL_{d_i}$ acts on the
    representation space $\operatorname{Rep}(Q, \mathbf{d}) = \bigoplus_{a \in Q_1}
    \Hom(\CC^{d_{s(a)}}, \CC^{d_{t(a)}})$ by change of basis;
  \item $W_{\mathbf{d}} \colon \operatorname{Rep}(Q, \mathbf{d}) \to \CC$ is
    the trace of the potential, $W_{\mathbf{d}}(\rho) = \Tr(W(\rho))$;
  \item $\varphi_{W_{\mathbf{d}}}$ is the sheaf of vanishing cycles of $W_{\mathbf{d}}$;
  \item $H^{\BM}_*$ denotes equivariant Borel--Moore homology.
\end{itemize}
\end{definition}

\begin{construction}[CoHA multiplication]
\label{constr:coha-mult}
The multiplication
\[
  \mu \colon \CoHA(Q,W)_{\mathbf{d}} \otimes \CoHA(Q,W)_{\mathbf{e}}
  \longrightarrow \CoHA(Q,W)_{\mathbf{d}+\mathbf{e}}
\]
is defined by the correspondence
\[
\begin{tikzcd}[column sep=large]
  & \left\{\begin{array}{c} 0 \to V' \to V \to V'' \to 0 \\ \text{short exact in } \operatorname{Rep}(Q) \end{array}\right\}
    \arrow[dl, "p"'] \arrow[dr, "q"] & \\
  \operatorname{Rep}(Q, \mathbf{d}) \times \operatorname{Rep}(Q, \mathbf{e})
  & & \operatorname{Rep}(Q, \mathbf{d}+\mathbf{e})
\end{tikzcd}
\]
where $p$ forgets to the sub and quotient and $q$ includes the extension.
Concretely, $\mu = q_! \circ p^*$, refined by the Thom--Sebastiani isomorphism
for vanishing cycles: $\varphi_{W_{\mathbf{d}}} \boxtimes \varphi_{W_{\mathbf{e}}}
\simeq \varphi_{W_{\mathbf{d}+\mathbf{e}}}|_{\text{extensions}}$.
\end{construction}

\begin{proposition}[Associativity]
\label{prop:coha-assoc}
The CoHA multiplication is associative.  That is, $(\CoHA(Q,W), \mu)$ is a
unital associative algebra; equivalently, an $E_1$-algebra in graded vector spaces.
\end{proposition}

\begin{proof}
Associativity follows from the transitivity of the extension correspondence:
a filtration $0 \subset V' \subset V'' \subset V$ with successive quotients
of dimensions $\mathbf{d}, \mathbf{e}, \mathbf{f}$ can be built either as
$(V' \subset V'') \subset V$ or $V' \subset (V'' \subset V)$, and the two
compositions yield the same pushforward.  The Thom--Sebastiani isomorphism
is itself associative, so the vanishing-cycle refinement preserves this.
\end{proof}

\begin{remark}[Noncommutativity]
\label{rem:noncomm}
The CoHA multiplication is \emph{not} commutative.  The extension correspondence
depends on which factor is the subrepresentation and which is the quotient;
this is the ``ordering'' that makes it $E_1$ rather than $E_2$.  The failure
of commutativity is measured precisely by the $R$-matrix (Section~\ref{subsec:r-matrix}).
\end{remark}


%% ======================================================================
\section{Precise statement: CoHA as $U(\frakn_+)$}
\label{sec:precise-statement}
%% ======================================================================

\subsection{Triangular decomposition of the quantum vertex chiral group}
\label{subsec:triangular}

\begin{definition}[Quantum vertex chiral group]
\label{def:qvcg}
For a CY3 $X$ with generalized root datum $\cR(X) = (\Lambda, \Delta^{\mathrm{re}},
\Delta^{\mathrm{im}}, \text{mult}, W, \rho)$ (see Chapter~1 of the main text),
the \emph{quantum vertex chiral group} $G(X)$ is the $E_2$-chiral algebra
$\Phi(X) = A_X$ whose:
\begin{enumerate}[label=(\roman*)]
  \item underlying Lie algebra $\frakg(X)$ is the generalized BKM superalgebra
    associated to $\cR(X)$;
  \item quantization is the affine (super) Yangian $Y(\widehat{\frakg}_{Q_X})$
    for toric $X$;
  \item representation category $\Rep^{E_2}(G(X))$ is braided monoidal with
    braiding given by the Yangian $R$-matrix.
\end{enumerate}
\end{definition}

\begin{definition}[Triangular decomposition]
\label{def:triangular-decomp}
The affine (super) Yangian $Y(\widehat{\frakg})$ admits a triangular decomposition
\[
  Y(\widehat{\frakg}) \cong Y^-(\widehat{\frakg}) \otimes Y^0(\widehat{\frakg})
    \otimes Y^+(\widehat{\frakg})
\]
where:
\begin{itemize}
  \item $Y^+(\widehat{\frakg})$ is generated by the positive modes $\{e_{i,r}\}_{i \in Q_0, r \geq 0}$;
  \item $Y^0(\widehat{\frakg})$ is the Cartan subalgebra generated by $\{\psi_{i,r}\}_{i \in Q_0, r \geq 0}$;
  \item $Y^-(\widehat{\frakg})$ is generated by the negative modes $\{f_{i,r}\}_{i \in Q_0, r \geq 0}$.
\end{itemize}
Each of $Y^{\pm}$ is an associative algebra (with no braiding);
the braiding arises from the interplay of $Y^+$ and $Y^-$ mediated by $Y^0$.
\end{definition}

\begin{theorem}[Critical CoHA $=$ positive half]
\label{thm:coha-positive-half}
Let $X$ be a toric CY3 without compact 4-cycles, with quiver-with-potential
$(Q_X, W_X)$.  There is an isomorphism of associative algebras
\[
  \boxed{\CoHA(Q_X, W_X) \;\cong\; Y^+(\widehat{\frakg}_{Q_X})}
\]
where $Y^+(\widehat{\frakg}_{Q_X})$ is the positive half of the affine super
Yangian associated to the toric quiver.  Under this isomorphism:
\begin{enumerate}[label=(\roman*)]
  \item The dimension-vector grading on $\CoHA$ corresponds to the root grading
    on $Y^+$.
  \item The CoHA generators in degree $\mathbf{e}_i$ (simple dimension vectors)
    correspond to the Yangian generators $e_{i,r}$.
  \item The CoHA relations (from geometry of extension correspondences) map to
    the Yangian Drinfeld relations.
  \item The equivariant parameters $(\epsilon_1, \epsilon_2, \epsilon_3)$ with
    $\epsilon_1 + \epsilon_2 + \epsilon_3 = 0$ correspond to the Yangian
    deformation parameters $(\hbar_1, \hbar_2, \hbar_3)$.
\end{enumerate}
\end{theorem}


\subsection{The universal enveloping algebra interpretation}
\label{subsec:uea}

The positive half $Y^+(\widehat{\frakg})$ should be thought of as a deformation
of $U(\frakn_+)$ where $\frakn_+$ is the positive nilpotent subalgebra of the
BKM superalgebra $\frakg(X)$:

\begin{proposition}
\label{prop:classical-limit}
In the classical limit $\hbar \to 0$ (non-equivariant limit $\epsilon_1 = \epsilon_2
= \epsilon_3 = 0$ after imposing $\epsilon_1 + \epsilon_2 + \epsilon_3 = 0$,
realized e.g.\ as $\epsilon_1 = \hbar, \epsilon_2 = -\hbar, \epsilon_3 = 0$
with $\hbar \to 0$), the positive Yangian degenerates:
\[
  Y^+(\widehat{\frakg}) \xrightarrow{\hbar \to 0} U(\frakn_+[\![z]\!])
\]
where $\frakn_+[\![z]\!]$ is the positive nilpotent of the loop algebra of $\frakg$.
The CoHA in the non-equivariant limit becomes the Hall algebra of the abelian
category of $\operatorname{Jac}(Q,W)$-modules, and the Hall algebra
$\to$ $U(\frakn_+)$ isomorphism is the classical Ringel theorem.
\end{proposition}


%% ======================================================================
\section{The $E_1$ structure: associativity and the ordered sector}
\label{sec:e1-structure}
%% ======================================================================

\subsection{CoHA as $E_1$-algebra}
\label{subsec:coha-e1}

The CoHA multiplication (Construction~\ref{constr:coha-mult}) makes $\CoHA(Q,W)$
into an $E_1$-algebra.  Let us unpack what $E_1$ means here concretely and
why the CoHA is \emph{not} naturally $E_2$.

\begin{proposition}[$E_1$-structure on CoHA]
\label{prop:coha-e1}
The CoHA $\CoHA(Q,W)$ is an $E_1$-algebra in the symmetric monoidal
$\infty$-category of graded mixed Hodge structures (or equivariant
$\ell$-adic sheaves).  Concretely:
\begin{enumerate}[label=(\roman*)]
  \item $\CoHA$ carries an associative multiplication $\mu$ (from
    short exact sequences);
  \item $\mu$ is \emph{not} commutative, not even up to coherent homotopy;
  \item the coherent associativity data (higher compositions from
    $n$-step filtrations for all $n \geq 2$) assembles into an
    $E_1$-algebra structure;
  \item there is no natural $E_2$-refinement of $\mu$ on $\CoHA$ alone.
\end{enumerate}
\end{proposition}

\begin{proof}
(i)--(iii): The $n$-ary multiplication $\mu_n$ on $\CoHA$ is defined by the
correspondence of $n$-step filtrations
\[
  0 = V_0 \subset V_1 \subset \cdots \subset V_n = V, \qquad
  \dim(V_k/V_{k-1}) = \mathbf{d}_k.
\]
The permutation group $S_n$ does \emph{not} act on this space (filtrations are
ordered), but the associahedron $K_n$ acts via refinement of filtrations.
This is the $E_1$-operad structure.

(iv): An $E_2$-structure would require a half-braiding $\sigma \colon \mu(a,b)
\xrightarrow{\sim} \mu(b,a)$ compatible with the $E_2$-operad.  The extension
correspondence is asymmetric: $\{0 \to V' \to V \to V'' \to 0\}$ is not
isomorphic to $\{0 \to V'' \to V \to V' \to 0\}$ as a correspondence.  The
discrepancy between these two orientations is precisely the $R$-matrix, which
requires passing to the Drinfeld double.
\end{proof}


\subsection{The Swiss-cheese interpretation}
\label{subsec:swiss-cheese}

The $E_1$-structure on the CoHA is the shadow of the Swiss-cheese algebra from
Volume~II.  The Swiss-cheese operad $\mathrm{SC}$ governs the
interaction of a bulk ($E_2$) algebra with a boundary ($E_1$) algebra.
In the CY3 context:

\begin{center}
\renewcommand{\arraystretch}{1.4}
\begin{tabular}{lll}
  \textbf{Swiss-cheese sector} & \textbf{Operadic structure} & \textbf{CY3 realization} \\
  \hline
  Bulk (interior of disk) & $E_2$ (braided) & Full Yangian $Y(\widehat{\frakg})$ \\
  Boundary (interval on $\partial$) & $E_1$ (associative) & CoHA $\cong Y^+(\widehat{\frakg})$ \\
  Bulk-boundary interaction & $\mathrm{SC}$-module & $Y^+$ is a module over $Y$
\end{tabular}
\end{center}

\noindent
The Swiss-cheese picture makes manifest why the CoHA is $E_1$: it lives on the
boundary of the $E_2$-disk.  The ``ordered'' direction is the real line forming
the boundary, while the full disk provides the $E_2$-enhancement.


\subsection{Geometric origin of associativity: tree-level curve counting}
\label{subsec:tree-level}

The CoHA multiplication encodes genus-0 (tree-level) enumerative geometry.
The correspondence
\[
  0 \to V' \to V \to V'' \to 0
\]
of short exact sequences in $\operatorname{Rep}(Q)$ is the quiver analogue of
the composition of stable maps:
\[
  \overline{\mathcal{M}}_{0,n_1+1}(X, \beta_1) \times_X
  \overline{\mathcal{M}}_{0,n_2+1}(X, \beta_2) \to
  \overline{\mathcal{M}}_{0,n_1+n_2}(X, \beta_1+\beta_2).
\]
Associativity of the CoHA corresponds to the associativity of the genus-0 OPE
(the WDVV equation in the GW language, or the quantum product on $QH^*(X)$).
This is an $E_1$-structure because tree-level amplitudes compose
\emph{sequentially} (the trees have a planar ordering from the root).


%% ======================================================================
\section{The $E_2$ structure: braiding and the $R$-matrix}
\label{sec:e2-structure}
%% ======================================================================

\subsection{The Drinfeld double}
\label{subsec:drinfeld-double}

The passage from the $E_1$-algebra $\CoHA$ to the $E_2$-algebra $G(X)$
proceeds via the Drinfeld double construction:

\begin{construction}[Drinfeld double of CoHA]
\label{constr:drinfeld-double}
Given the CoHA $\CoHA = Y^+(\widehat{\frakg})$, define the dual CoHA
\[
  \CoHA^{\vee} = Y^-(\widehat{\frakg})
\]
as the CoHA of the opposite quiver $(Q^{\op}, W^{\op})$ (reverse all arrows,
negate the potential).  The Drinfeld double is the algebra generated by $\CoHA$
and $\CoHA^{\vee}$ together with a Cartan subalgebra $Y^0$, subject to the
cross-relations
\[
  [e_{i,r},\, f_{j,s}] = \delta_{ij}\, \psi_{i, r+s}
\]
and the Yangian exchange relations between $Y^0$ and $Y^{\pm}$.  The result
is the full Yangian:
\[
  \operatorname{Drin}(\CoHA(Q,W)) \cong Y(\widehat{\frakg}_{Q}).
\]
\end{construction}

\begin{theorem}[Drinfeld double $=$ full Yangian]
\label{thm:drinfeld-double}
The Drinfeld double of $\CoHA(Q,W)$ (with respect to the topological bialgebra
structure induced by the coproduct dual to the CoHA multiplication on
$\CoHA^{\vee}$) is isomorphic to the full affine super Yangian:
\[
  \operatorname{Drin}(\CoHA(Q,W)) \cong Y(\widehat{\frakg}_{Q}).
\]
The universal $R$-matrix of this Drinfeld double satisfies the quantum
Yang--Baxter equation and provides the $E_2$-braiding.
\end{theorem}


\subsection{The $R$-matrix and the passage $E_1 \to E_2$}
\label{subsec:r-matrix}

\begin{construction}[$R$-matrix from the Drinfeld double]
\label{constr:r-matrix}
The universal $R$-matrix of $Y(\widehat{\frakg})$ is the element
\[
  R \in Y^+(\widehat{\frakg}) \;\widehat{\otimes}\;
      Y^-(\widehat{\frakg}) \subset
      Y(\widehat{\frakg}) \;\widehat{\otimes}\;
      Y(\widehat{\frakg})
\]
characterized (up to scalar) by the intertwining property
\[
  R \cdot \Delta(x) = \Delta^{\op}(x) \cdot R \qquad
  \text{for all } x \in Y(\widehat{\frakg})
\]
where $\Delta$ is the coproduct and $\Delta^{\op} = \sigma \circ \Delta$ is
the opposite coproduct.  The $R$-matrix satisfies the quantum Yang--Baxter equation
\[
  R_{12} R_{13} R_{23} = R_{23} R_{13} R_{12}
  \quad \in Y(\widehat{\frakg})^{\widehat{\otimes}\, 3}.
\]
\end{construction}

\begin{proposition}[$R$-matrix provides $E_2$-structure]
\label{prop:r-matrix-e2}
The data $(\CoHA, \CoHA^{\vee}, Y^0, R)$ assemble into an $E_2$-algebra structure
on the full Yangian $Y(\widehat{\frakg})$:
\begin{enumerate}[label=(\roman*)]
  \item The first $E_1$-direction comes from the CoHA multiplication
    (positive modes, extension of representations);
  \item The second $E_1$-direction comes from the Cartan/coweight lattice
    (dimension-vector grading and the $\psi$-generators);
  \item By Dunn additivity $E_2 \simeq E_1 \otimes E_1$, these two
    $E_1$-structures combine into an $E_2$-structure;
  \item The braiding $\sigma \colon V \otimes W \xrightarrow{R}
    W \otimes V$ in $\Rep^{E_2}(Y(\widehat{\frakg}))$ is given by the
    action of $R$.
\end{enumerate}
\end{proposition}

\begin{proof}
The key point is the Dunn additivity decomposition.  An $E_2$-algebra is an
$E_1$-algebra \emph{in the category of $E_1$-algebras}.  For the Yangian:
\begin{itemize}
  \item The ``outer'' $E_1$-structure: the CoHA multiplication (extending
    representations, increasing dimension vector).  This is associative and
    makes $Y^+$ into an $E_1$-algebra.
  \item The ``inner'' $E_1$-structure: the mode grading / spectral-parameter
    direction.  The generators $\{e_{i,r}\}_{r \geq 0}$ at fixed vertex $i$
    form an $E_1$-algebra (the modes of a vertex operator).
  \item The compatibility of these two $E_1$-structures (the Drinfeld
    exchange relations between the dimension-vector direction and the
    spectral-parameter direction) is precisely the data of an $E_2$-structure.
\end{itemize}
The $R$-matrix encodes the non-trivial braiding that distinguishes $E_2$
from $E_{\infty}$; it measures the failure of commutativity of the two
$E_1$-directions.
\end{proof}


\subsection{The chiral algebra perspective}
\label{subsec:chiral-perspective}

In the language of the main text (Chapters 5--6), the quantum vertex chiral
group $G(X)$ is an $E_2$-chiral algebra on $\Ran(X) \times \Ran(Y)$ with:

\begin{itemize}
  \item The $X$-direction (holomorphic) encodes the OPE (vertex algebra
    structure, spectral parameter $z$).
  \item The $Y$-direction (holomorphic) encodes the braiding (the second
    ``quantum'' direction, spectral parameter $w$).
  \item The $R$-matrix $R(z-w)$ is the two-point function in the $Y$-direction.
  \item Restricting to a single point in the $Y$-direction (i.e., forgetting
    the braiding) recovers the $E_1$-chiral algebra $= \CoHA$.
\end{itemize}

\noindent
The $E_2$ bar complex $B_{E_2}(G(X))$ (see the main text, Chapter 9) has:
\begin{itemize}
  \item Two commuting differentials $d_X, d_Y$ from the two holomorphic
    directions;
  \item Two commuting coproducts $\Delta_X, \Delta_Y$;
  \item The $R$-matrix as the transposition datum between the two
    $E_1$-sectors.
\end{itemize}


%% ======================================================================
\section{Proof for $\CC^3$: Schiffmann--Vasserot}
\label{sec:proof-c3}
%% ======================================================================

We now verify Theorem~\ref{thm:coha-positive-half} in complete detail for
the simplest case $X = \CC^3$.

\subsection{Setup}
\label{subsec:c3-setup}

For $\CC^3$, the quiver-with-potential is:
\begin{itemize}
  \item $Q$: one vertex, three loops $x, y, z$ (the ``tripled Jordan quiver'');
  \item $W = \Tr(x[y,z]) = \Tr(xyz - xzy)$.
\end{itemize}
The representation space of dimension $d$ is
$\operatorname{Rep}(Q, d) = \operatorname{Mat}_{d \times d}(\CC)^3$
with $\GL_d$-action by simultaneous conjugation.  The function
$W_d \colon \operatorname{Mat}_d^3 \to \CC$ is
$W_d(X,Y,Z) = \Tr(X[Y,Z])$.


\subsection{The CoHA for $\CC^3$}
\label{subsec:c3-coha}

The equivariant CoHA is
\[
  \CoHA(\CC^3) = \bigoplus_{d \geq 0}
    H^{\BM,\, T}_*(\Crit(W_d) / \GL_d,\, \varphi_{W_d})
\]
where $T = (\CC^*)^3$ acts on $\CC^3$ with weights $\epsilon_1, \epsilon_2,
\epsilon_3$, inducing an action on $\operatorname{Mat}_d^3$ where $X, Y, Z$
have weights $\epsilon_1, \epsilon_2, \epsilon_3$ respectively.

\begin{proposition}[Kontsevich--Soibelman, Schiffmann--Vasserot]
\label{prop:c3-coha-generators}
As a $\QQ[\epsilon_1, \epsilon_2, \epsilon_3]/(\epsilon_1+\epsilon_2+\epsilon_3)$-algebra,
$\CoHA(\CC^3)$ is generated by elements
\[
  e_r \in \CoHA(\CC^3)_{d=1}, \qquad r = 0, 1, 2, \ldots
\]
corresponding to the equivariant classes
$e_r = [pt]^{\BM} \cdot u^r \in H^{\BM,T}_*(\CC^3/\CC^*) \cong \QQ[u]$
where $u$ is the equivariant parameter of the $\GL_1 = \CC^*$ framing.
\end{proposition}


\subsection{The affine Yangian $Y(\widehat{\gl}_1)$}
\label{subsec:affine-yangian-gl1}

\begin{definition}[Affine Yangian of $\gl_1$]
\label{def:affine-yangian-gl1}
The affine Yangian $Y(\widehat{\gl}_1)$ (also denoted $Y_{\hbar_1, \hbar_2, \hbar_3}$
with $\hbar_1 + \hbar_2 + \hbar_3 = 0$) is the associative algebra with
generators
\[
  e_r, \quad f_r, \quad \psi_r, \qquad r = 0, 1, 2, \ldots
\]
and relations:
\begin{align}
  [\psi_r, \psi_s] &= 0, \label{eq:yangian-cartan} \\
  [\psi_0, e_r] &= 0, \quad [\psi_0, f_r] = 0, \label{eq:yangian-psi0} \\
  [\psi_{r+1}, e_s] - [\psi_r, e_{s+1}] &= \sigma_2\, \{\psi_r, e_s\}, \label{eq:yangian-psi-e} \\
  [\psi_{r+1}, f_s] - [\psi_r, f_{s+1}] &= -\sigma_2\, \{\psi_r, f_s\}, \label{eq:yangian-psi-f} \\
  [e_{r+1}, e_s] - [e_r, e_{s+1}] &= \sigma_2\, \{e_r, e_s\}, \label{eq:yangian-ee} \\
  [f_{r+1}, f_s] - [f_r, f_{s+1}] &= -\sigma_2\, \{f_r, f_s\}, \label{eq:yangian-ff} \\
  [e_r, f_s] &= \psi_{r+s}, \label{eq:yangian-ef} \\
  \operatorname{Sym}_{r_1, r_2, r_3}\, [e_{r_1}, [e_{r_2}, e_{r_3}]] &= 0, \label{eq:yangian-serre-e} \\
  \operatorname{Sym}_{r_1, r_2, r_3}\, [f_{r_1}, [f_{r_2}, f_{r_3}]] &= 0, \label{eq:yangian-serre-f}
\end{align}
where $\sigma_2 = \hbar_1 \hbar_2 + \hbar_1 \hbar_3 + \hbar_2 \hbar_3$ and
$\{a,b\} = ab + ba$ is the anticommutator.
\end{definition}

\begin{remark}[Parameters]
\label{rem:parameters}
The three parameters $\hbar_1, \hbar_2, \hbar_3$ with
$\hbar_1 + \hbar_2 + \hbar_3 = 0$ correspond to the equivariant parameters
$\epsilon_1, \epsilon_2, \epsilon_3$ of the torus acting on $\CC^3$.  The
constraint $\epsilon_1 + \epsilon_2 + \epsilon_3 = 0$ is the CY condition:
the canonical bundle $\omega_{\CC^3} \cong \cO_{\CC^3}$ has weight
$\epsilon_1 + \epsilon_2 + \epsilon_3 = 0$.
\end{remark}

\begin{theorem}[Schiffmann--Vasserot \cite{SV-elliptic-hall, SV-cherednik-aha}]
\label{thm:sv}
There is an isomorphism of associative algebras
\[
  \CoHA(\CC^3) \;\xrightarrow{\;\sim\;}\; Y^+(\widehat{\gl}_1)
\]
sending the CoHA generators $e_r \in \CoHA_{d=1}$ to the Yangian
generators $e_r$.  The CoHA multiplication corresponds to the Yangian
relations \eqref{eq:yangian-ee}--\eqref{eq:yangian-serre-e}.
\end{theorem}

\begin{proof}[Proof outline]
The proof proceeds in three steps.

\emph{Step 1: Identify generators.}
The dimension-$1$ piece $\CoHA_{d=1} = H^{\BM,T}_*(\CC^3/\CC^*,\, \varphi_{W_1})$
is the localized equivariant Borel--Moore homology of a point (since $W_1 = 0$
for $d=1$ and $\Crit(0) = \CC^3$).  After $T$-equivariant localization,
$\CoHA_{d=1} \cong \QQ[u]$ with generator $e_r = u^r$.

\emph{Step 2: Compute the multiplication.}
The CoHA multiplication $\mu \colon \CoHA_{d=1}^{\otimes 2} \to \CoHA_{d=2}$
is computed by the correspondence of short exact sequences
$0 \to \CC \to \CC^2 \to \CC \to 0$ of quiver representations.
By equivariant localization on the flag variety
$\operatorname{Fl}(1,2; \CC^2) = \mathbb{P}^1$, one computes
\[
  [e_{r+1}, e_s] - [e_r, e_{s+1}] = \sigma_2\, \{e_r, e_s\}
\]
where $\sigma_2 = \epsilon_1\epsilon_2 + \epsilon_1\epsilon_3 + \epsilon_2\epsilon_3$
arises from the weights of the normal bundle to the fixed locus.

\emph{Step 3: Verify the Serre relation.}
The cubic Serre relation $\operatorname{Sym}_{r_1,r_2,r_3} [e_{r_1},[e_{r_2},e_{r_3}]] = 0$
follows from a computation in $\CoHA_{d=3}$.  The symmetric group $S_3$ acts
on the space of 3-step filtrations, and the vanishing of the symmetrization
follows from the vanishing of the relevant equivariant class on the full flag
variety $\operatorname{Fl}(1,2,3;\CC^3)$.  This is a consequence of
the identity $\sum_{\sigma \in S_3} \operatorname{sgn}(\sigma) \prod_{i<j}
(\sigma(u_i) - \sigma(u_j) + \sigma_2) = 0$ in the localized equivariant
cohomology.
\end{proof}


\subsection{The Drinfeld double: $Y(\widehat{\gl}_1) \cong \cW_{1+\infty}$}
\label{subsec:drinfeld-double-c3}

\begin{theorem}[Schiffmann--Vasserot, Maulik--Okounkov, Tsymbaliuk]
\label{thm:drinfeld-double-c3}
The Drinfeld double of the $\CC^3$ CoHA is the full affine Yangian:
\[
  \operatorname{Drin}(\CoHA(\CC^3)) \cong Y(\widehat{\gl}_1).
\]
The algebra $Y(\widehat{\gl}_1)$ is isomorphic to the universal enveloping
algebra of the $\cW_{1+\infty}$-algebra at the self-dual level
$c = \hbar_1/\hbar_2 + \hbar_2/\hbar_1 + 2$:
\[
  Y(\widehat{\gl}_1) \cong U(\cW_{1+\infty}).
\]
\end{theorem}

\begin{proof}[Proof outline]
The identification $Y(\widehat{\gl}_1) \cong U(\cW_{1+\infty})$ was established
by Schiffmann--Vasserot via the spherical Hall algebra of the elliptic curve
(as a limit of the elliptic Hall algebra / Ding--Iohara--Miki algebra), and
independently by Maulik--Okounkov via stable envelopes and $R$-matrix
constructions.

The key steps are:
\begin{enumerate}[label=(\arabic*)]
  \item The CoHA $\CoHA(\CC^3) = Y^+(\widehat{\gl}_1)$ is the algebra of
    creation operators on the Fock space
    $\cF = \bigoplus_{d \geq 0} H^{\BM,T}_*(\operatorname{Hilb}^d(\CC^3))$.
  \item The dual CoHA $\CoHA^{\vee}(\CC^3) = Y^-(\widehat{\gl}_1)$ consists
    of annihilation operators.
  \item The Cartan subalgebra $Y^0(\widehat{\gl}_1)$ is generated by the
    operators $\psi_r$ which act diagonally on the fixed-point basis
    $\{|\lambda\rangle\}_{\lambda \in \mathcal{P}}$ indexed by 3D partitions.
  \item The cross-relation $[e_r, f_s] = \psi_{r+s}$ follows from the
    geometric identity expressing the commutator of creation/annihilation
    as a diagonal operator (the ``boson-fermion'' computation on the Fock space).
  \item The isomorphism with $\cW_{1+\infty}$ sends $e_0 \mapsto W_1 = J$
    (the $U(1)$ current), $e_1 \mapsto W_2 = T$ (the Virasoro field),
    and the higher $e_r$ to the higher-spin $W$-algebra generators $W_{r+1}$.
\end{enumerate}
\end{proof}

\begin{remark}[Connection to Volume I]
\label{rem:vol1-connection}
The identification $Y(\widehat{\gl}_1) \cong \cW_{1+\infty}$ is precisely
the statement that the quantum vertex chiral group of $\CC^3$ is the
$\cW_{1+\infty}$-algebra, the object whose modular Koszul duality was
studied extensively in Volume~I (the $W$-infinity deep audit, Session~8).
The modular characteristic $\kappa(\cW_{1+\infty})$ controls the genus-$g$
DT invariants of $\CC^3$ via the obstruction $\mathrm{obs}_g = \kappa \cdot \lambda_g$.
\end{remark}


\subsection{The $R$-matrix for $\CC^3$}
\label{subsec:r-matrix-c3}

\begin{construction}[Maulik--Okounkov $R$-matrix]
\label{constr:mo-r-matrix}
The Maulik--Okounkov stable envelope construction produces an explicit
$R$-matrix
\[
  R(z) \in \End(\cF \otimes \cF)[\![z^{-1}]\!]
\]
acting on the tensor product of two Fock spaces.  This $R$-matrix satisfies:
\begin{enumerate}[label=(\roman*)]
  \item The quantum Yang--Baxter equation
    $R_{12}(z_1-z_2) R_{13}(z_1-z_3) R_{23}(z_2-z_3) =
     R_{23}(z_2-z_3) R_{13}(z_1-z_3) R_{12}(z_1-z_2)$.
  \item $R(z)$ intertwines $\Delta$ and $\Delta^{\op}$:
    $R(z) \Delta(x) = \Delta^{\op}(x) R(z)$.
  \item In the classical limit $\epsilon_1, \epsilon_2 \to 0$,
    $R(z) = 1 + r(z)/z + O(z^{-2})$ where $r(z)$ is the classical $r$-matrix
    of $\widehat{\gl}_1$.
  \item In the chiral algebra language of the main text,
    $R(z) = \mathrm{Res}^{\mathrm{coll}}_{0,2}(\Theta_{A_{\CC^3}})$
    is the collision $R$-matrix extracted from the universal MC element.
\end{enumerate}
\end{construction}

\begin{proposition}[$E_1 \to E_2$ for $\CC^3$]
\label{prop:e1-to-e2-c3}
The passage from the CoHA ($E_1$) to the full $\cW_{1+\infty}$ ($E_2$) for $\CC^3$
is given by:
\[
\begin{tikzcd}[column sep=huge]
  \CoHA(\CC^3) = Y^+(\widehat{\gl}_1)
    \arrow[r, hook, "\text{Drin. double}"]
  & Y(\widehat{\gl}_1) \cong \cW_{1+\infty}
    \arrow[r, "\text{$R$-matrix}"]
  & \Rep^{E_2}(Y(\widehat{\gl}_1)) \text{ braided}
\end{tikzcd}
\]
where:
\begin{itemize}
  \item the first arrow is the inclusion $Y^+ \hookrightarrow
    \operatorname{Drin}(Y^+) = Y$;
  \item the second arrow equips the representation category with the
    braiding $\sigma_{V,W} = P \circ R(z)$ (permutation composed with
    $R$-matrix);
  \item the braided monoidal category $\Rep^{E_2}(Y(\widehat{\gl}_1))$
    is the $E_2$-representation category of the quantum vertex chiral group
    $G(\CC^3)$.
\end{itemize}
\end{proposition}


%% ======================================================================
\section{General toric CY3: the RSYZ theorem}
\label{sec:rsyz}
%% ======================================================================

\subsection{Toric data and quivers}
\label{subsec:toric-data}

A toric CY3 $X_\Sigma$ is determined by a fan $\Sigma$ in $\ZZ^3$ with all
ray generators on a common affine hyperplane (the CY condition: $c_1(X) = 0$
in the toric sense).  The toric diagram is the convex hull of the projected
generators.

\begin{construction}[Quiver from toric diagram]
\label{constr:quiver-toric}
The brane tiling / dimer model construction associates to a toric diagram
$\Delta$ a quiver with potential $(Q_X, W_X)$:
\begin{enumerate}[label=(\roman*)]
  \item Faces of the brane tiling $\to$ vertices of $Q_X$;
  \item Edges of the brane tiling $\to$ arrows of $Q_X$;
  \item Elementary cycles of the tiling $\to$ terms of $W_X$.
\end{enumerate}
The resulting $(Q_X, W_X)$ is such that $\operatorname{Jac}(Q_X, W_X)$ is the
coordinate ring of $X$ (as a 3-CY algebra).
\end{construction}

\begin{example}[Classification of simple toric CY3]
\label{ex:toric-classification}
\leavevmode
\begin{enumerate}[label=(\alph*)]
  \item $\CC^3$: triangle diagram, Jordan quiver, $W = \Tr(x[y,z])$.
  \item Resolved conifold $\cO(-1) \oplus \cO(-1) \to \mathbb{P}^1$:
    two-triangle diagram, Klebanov--Witten quiver.
  \item $\CC^3/\ZZ_N$ (orbifolds): $N$-gon diagram, McKay quiver of
    $\ZZ_N \subset SL_3(\CC)$.
  \item $(\CC^*)^3$ (open torus): infinite periodic tiling, the ``affine''
    case with infinitely many dimension vectors.
\end{enumerate}
\end{example}

\subsection{The RSYZ theorem}
\label{subsec:rsyz-thm}

\begin{theorem}[Rapcak--Soibelman--Yang--Zhao]
\label{thm:rsyz}
Let $X$ be a toric CY3 without compact 4-cycles.  Then:
\begin{enumerate}[label=(\roman*)]
  \item There exists an affine super Yangian $Y(\widehat{\frakg}_{Q_X})$
    associated to the toric quiver $Q_X$, with the super Lie algebra
    $\frakg_{Q_X}$ determined by the toric data.
  \item The critical CoHA is isomorphic to the positive half:
    \[
      \CoHA(Q_X, W_X) \cong Y^+(\widehat{\frakg}_{Q_X}).
    \]
  \item The full affine super Yangian $Y(\widehat{\frakg}_{Q_X})$ is the
    Drinfeld double of the CoHA.
  \item The representation category $\Rep(Y(\widehat{\frakg}_{Q_X}))$
    acts on the equivariant BPS cohomology
    $\bigoplus_{\beta, n} H^{\BM,T}_*(\cM_X(\beta, n))$ of the moduli
    of sheaves on $X$.
\end{enumerate}
\end{theorem}

\begin{remark}[The ``no compact 4-cycles'' condition]
\label{rem:no-compact-4}
The condition that $X$ has no compact 4-cycles (equivalently, the toric
diagram is convex with no interior lattice points other than those on
the boundary) ensures that the CoHA has a well-defined associative
multiplication.  When compact 4-cycles are present, additional
``wall-crossing'' phenomena occur and the CoHA structure is modified
by quantum corrections.
\end{remark}


\subsection{The generalized root datum from toric geometry}
\label{subsec:root-datum}

\begin{construction}[Root datum $\cR(X_\Sigma)$]
\label{constr:root-datum-toric}
For a toric CY3 $X_\Sigma$, the generalized root datum of $G(X)$:
\begin{enumerate}[label=(\roman*)]
  \item \textbf{Root lattice}: $\Lambda(X) = \ZZ^{Q_0}$ (dimension vectors),
    with Euler form $\chi(\mathbf{d}, \mathbf{e}) =
    \sum_{i \in Q_0} d_i e_i - \sum_{a \in Q_1} d_{s(a)} e_{t(a)}$
    and symmetrized Euler form
    $(\mathbf{d}, \mathbf{e}) = \chi(\mathbf{d}, \mathbf{e}) + \chi(\mathbf{e}, \mathbf{d})$.
  \item \textbf{Real simple roots}: the simple dimension vectors
    $\{\mathbf{e}_i\}_{i \in Q_0}$, with parity $|\mathbf{e}_i| \in \ZZ/2$ determined
    by the number of loops at vertex $i$: $|\mathbf{e}_i| = 0$ if the number
    of loops is even, $|\mathbf{e}_i| = 1$ if odd.
  \item \textbf{Generalized Cartan matrix}: $a_{ij} = (\mathbf{e}_i, \mathbf{e}_j)$,
    the symmetrized Euler form evaluated on simple roots.
  \item \textbf{Imaginary roots}: dimension vectors $\mathbf{d}$ with
    $\DT_{\mathbf{d}}(X) \neq 0$, with multiplicity
    $\operatorname{mult}(\mathbf{d}) = |\DT_{\mathbf{d}}(X)|$.
  \item \textbf{Weyl group}: $W(X)$ generated by the reflections
    $s_i(\mathbf{d}) = \mathbf{d} - a_{ii}^{-1}(\mathbf{d}, \mathbf{e}_i)\, \mathbf{e}_i$
    for real simple roots.
\end{enumerate}
This is the data needed to define the BKM superalgebra $\frakg_{Q_X}$
and its affinization $\widehat{\frakg}_{Q_X}$.
\end{construction}

\begin{proposition}[Toric fan $=$ Dynkin data]
\label{prop:fan-dynkin}
The generalized Cartan matrix $a_{ij} = (\mathbf{e}_i, \mathbf{e}_j)$ is
determined by the toric diagram: the number of arrows between vertices $i$ and $j$
in $Q_X$ equals the number of edges shared by the corresponding faces in the brane
tiling.  In particular:
\begin{itemize}
  \item $a_{ii} = 2 - 2\ell_i$ where $\ell_i$ is the number of loops at vertex $i$;
  \item $a_{ij} = -($ number of arrows from $i$ to $j$ $+$ number of arrows from
    $j$ to $i)$ for $i \neq j$.
\end{itemize}
The super structure (even/odd roots) is determined by the diagonal: $\mathbf{e}_i$
is odd (fermionic) if and only if $a_{ii} < 0$, i.e., the vertex $i$ has
more than one loop.  For $\CC^3$: the single vertex has three loops, so
$a_{11} = 2 - 6 = -4$, the unique simple root is odd, and $\frakg_{Q_{\CC^3}}$
is a super Lie algebra.
\end{proposition}


%% ======================================================================
\section{The $E_1$/$E_2$ dictionary}
\label{sec:e1-e2-dictionary}
%% ======================================================================

We summarize the full $E_1 \to E_2$ passage as a dictionary:

\begin{center}
\renewcommand{\arraystretch}{1.5}
\begin{tabular}{lll}
  \textbf{$E_1$-sector (CoHA)} & \textbf{$E_2$-full (Yangian)} & \textbf{Passage} \\
  \hline
  $\CoHA(Q,W)$ & $Y(\widehat{\frakg}_Q)$ & Drinfeld double \\
  Associative multiplication $\mu$ & Braided product $\mu_R$ & $R$-matrix \\
  $\Rep^{E_1}(\CoHA)$ monoidal & $\Rep^{E_2}(Y)$ braided monoidal & Drinfeld center \\
  Short exact sequences & Yang--Baxter equation & Wall-crossing \\
  Genus-0 DT / tree-level & All-genus DT / all-loop & Loop expansion \\
  $Y^+(\widehat{\frakg})$ & $Y^+ \otimes Y^0 \otimes Y^-$ & Triangular decomp. \\
  One holomorphic direction & Two holomorphic directions & Second $E_1$-dir. \\
  Ordered OPE (planar) & Braided OPE (non-planar) & Braiding \\
  $\frakn_+$ (nilpotent) & $\frakg = \frakn_- \oplus \frakh \oplus \frakn_+$ & Borel--Weyl \\
  Creation operators on $\cF$ & Full $\cW$-algebra & Add annihilation \\
\end{tabular}
\end{center}

\begin{remark}[The Drinfeld center perspective]
\label{rem:drinfeld-center}
The passage from $E_1$ to $E_2$ has a clean categorical description via the
Drinfeld center (Chapter~12 of the main text).  The monoidal category
$\Rep^{E_1}(\CoHA)$ is a monoidal category.  Its Drinfeld center
\[
  \cZ(\Rep^{E_1}(\CoHA)) \simeq \Rep^{E_2}(Y(\widehat{\frakg}))
\]
is the braided monoidal representation category of the full Yangian.  This is
the categorical incarnation of the Drinfeld double: the center construction
takes an $E_1$-monoidal category and produces an $E_2$-braided monoidal category.
In the bulk-boundary language:
\begin{itemize}
  \item $E_1$ (boundary): the CoHA lives on the boundary of a 2-disk,
    corresponding to the ordered composition of representations;
  \item $E_2$ (bulk): the full Yangian lives in the interior of the 2-disk,
    where representations can braid around each other;
  \item The center construction is the ``bulk-boundary correspondence''
    $Z^{\mathrm{der}}_{\mathrm{ch}}(\CoHA) = Y(\widehat{\frakg})$.
\end{itemize}
\end{remark}


%% ======================================================================
\section{Refinements and generalizations}
\label{sec:refinements}
%% ======================================================================

\subsection{Cohomological refinement}
\label{subsec:coh-refinement}

The CoHA carries a natural cohomological (perverse) filtration from the
mixed Hodge structure on the Borel--Moore homology of vanishing cycles.
The associated graded with respect to this filtration is the BPS cohomology:

\begin{proposition}[BPS cohomology]
\label{prop:bps}
There is a ``PBW filtration'' on the CoHA such that
\[
  \operatorname{gr}^{\mathrm{PBW}} \CoHA(Q,W)
  \cong \Sym\Bigl(\bigoplus_{\mathbf{d}} \operatorname{BPS}_{\mathbf{d}}(Q,W) \otimes
    \QQ[u_1, \ldots, u_{d_1}, \ldots]\Bigr)
\]
where $\operatorname{BPS}_{\mathbf{d}}$ are the ``BPS invariants,'' the primitive
generators of the CoHA.  These are the root spaces of the BKM superalgebra:
$\operatorname{BPS}_{\mathbf{d}} = \frakg(X)_{\mathbf{d}}$ (the weight-$\mathbf{d}$
root space).
\end{proposition}


\subsection{K-theoretic and elliptic lifts}
\label{subsec:k-theoretic}

The cohomological CoHA has K-theoretic and elliptic upgrades:

\begin{center}
\renewcommand{\arraystretch}{1.4}
\begin{tabular}{llll}
  \textbf{Theory} & \textbf{CoHA} & \textbf{Quantum group} & \textbf{$E_2$-algebra} \\
  \hline
  Cohomological & $H^{\BM}_*(\Crit, \varphi)$ & Yangian $Y(\widehat{\frakg})$ & Affine Yangian \\
  K-theoretic & $K_T(\Crit, \varphi)$ & Quantum toroidal $U_{q,t}(\ddot{\frakg})$ & Quantum toroidal \\
  Elliptic & $\operatorname{Ell}_T(\Crit, \varphi)$ & Elliptic quantum group & Elliptic QG \\
\end{tabular}
\end{center}

\noindent
Each row has the same $E_1$/$E_2$ structure: the CoHA (positive half, $E_1$)
embeds via Drinfeld double into the full quantum group ($E_2$).  The passage
from cohomology to K-theory to elliptic cohomology corresponds to the passage
from Yangian to quantum toroidal to elliptic quantum group.

For $\CC^3$:
\begin{itemize}
  \item Cohomological: $Y(\widehat{\gl}_1) \cong \cW_{1+\infty}$;
  \item K-theoretic: the quantum toroidal algebra $U_{q,t}(\ddot{\gl}_1)$,
    also known as the Ding--Iohara--Miki algebra;
  \item Elliptic: the elliptic Hall algebra of Burban--Schiffmann, which is
    the original form in which the $\CC^3$ story was discovered.
\end{itemize}


\subsection{Wall-crossing and the quantum dilogarithm}
\label{subsec:wall-crossing}

The quantum dilogarithm identity underlying DT wall-crossing has a clean
interpretation in the $E_1$/$E_2$ framework:

\begin{proposition}[Wall-crossing as $E_1$/$E_2$ exchange]
\label{prop:wall-crossing}
The Kontsevich--Soibelman wall-crossing formula
\[
  \prod_{\gamma \in C_+}^{\curvearrowleft}
  \mathbb{E}(\DT_\gamma \cdot \hat{y}^\gamma)
  =
  \prod_{\gamma \in C_-}^{\curvearrowleft}
  \mathbb{E}(\DT'_\gamma \cdot \hat{y}^\gamma)
\]
(where $\mathbb{E}$ is the quantum dilogarithm and $\curvearrowleft$ denotes
the clockwise-ordered product) expresses the invariance of the $E_2$-structure
under change of $E_1$-splitting.  Concretely:
\begin{itemize}
  \item Different choices of stability condition $\zeta$ determine different
    ``positive halves'' $\CoHA_\zeta$ (different $E_1$-subalgebras of the
    $E_2$-algebra).
  \item The wall-crossing formula relates the ordered products (in different
    $E_1$-structures) via the $R$-matrix.
  \item The quantum dilogarithm $\mathbb{E}$ is the generating function of
    the CoHA in each chamber.
\end{itemize}
\end{proposition}


\subsection{Beyond toric: general CY3}
\label{subsec:beyond-toric}

For non-toric CY3 (compact CY3, CY3 with non-abelian gauge theory), the
CoHA is still defined but the identification with a Yangian is more subtle:

\begin{conjecture}[General CY3 CoHA conjecture]
\label{conj:general-cy3}
For any CY3 $X$ with a good moduli theory (e.g., a projective CY3 with an
appropriate stability condition), the critical CoHA $\CoHA(X)$ is the positive
half of a quantum vertex chiral group $G(X)$:
\[
  \CoHA(X) \cong U(\frakn_+(X))
\]
where $\frakn_+(X)$ is a pro-nilpotent Lie algebra whose root spaces are the
BPS cohomologies of $X$.  The Drinfeld double $\operatorname{Drin}(\CoHA(X))$
is an algebra whose representation category is the braided monoidal category
governing the BPS structure of $X$.
\end{conjecture}

\noindent
Evidence for this conjecture:
\begin{enumerate}[label=(\alph*)]
  \item For toric CY3: Theorem~\ref{thm:rsyz} (proved).
  \item For local curves $X = \operatorname{Tot}(\omega_C)$: the CoHA is
    related to the Yangian of $\gl_1$ (or $\gl_r$ for higher rank), and the
    Drinfeld double is the quantum group governing Pandharipande--Thomas /
    Gopakumar--Vafa invariants.
  \item For compact CY3: the Joyce--Song / Kontsevich--Soibelman wall-crossing
    framework implicitly constructs the CoHA and its Drinfeld double, although
    the explicit Yangian identification remains open.
  \item For the $K3 \times E$ family (Chapter~14 of the main text): the
    quantum vertex chiral group is the BKM superalgebra $\frakg_{\Delta_5}$,
    and the CoHA should be $U(\frakn_+(\frakg_{\Delta_5}))$.
\end{enumerate}


%% ======================================================================
\section{Summary diagram}
\label{sec:summary}
%% ======================================================================

The full structure can be summarized in the following diagram:

\[
\begin{tikzcd}[column sep=small, row sep=large]
  & \text{CY3 } X \arrow[dl, "\text{quiver}"'] \arrow[dr, "\text{root datum}"] & \\
  (Q_X, W_X) \arrow[d, "\text{BM homology}"] & &
    \cR(X) \arrow[d, "\text{BKM algebra}"] \\
  \CoHA(Q_X, W_X) \arrow[r, equal, "\text{Thm \ref{thm:coha-positive-half}}"]
    \arrow[d, "\text{Drin. double}"]
  & Y^+(\widehat{\frakg}_{Q_X}) \arrow[r, hook]
  & \frakg_{Q_X} \\
  Y(\widehat{\frakg}_{Q_X}) \arrow[r, "\text{$R$-matrix}"]
  & \Rep^{E_2}(Y(\widehat{\frakg}_{Q_X}))
    \arrow[r, equal, "\text{Thm CY-C}"]
  & \Rep^{E_2}(G(X))
\end{tikzcd}
\]

\noindent
Reading down the left column: CY3 geometry $\to$ CoHA ($E_1$) $\to$ full Yangian
($E_2$).  The critical CoHA is the $E_1$-sector (the positive half, the
ordered part, the creation operators, the tree-level) of the quantum vertex
chiral group.  The full $E_2$-structure adds the braiding, the $R$-matrix, the
Cartan and negative parts, and the all-genus corrections.

%% ======================================================================
%% References (informal, for a working note)
%% ======================================================================

\begin{thebibliography}{99}

\bibitem{SV-elliptic-hall}
O.~Schiffmann and E.~Vasserot,
\emph{The elliptic Hall algebra and the $K$-theory of the Hilbert scheme of $\mathbb{A}^2$},
Duke Math.\ J.\ \textbf{162} (2013), no.~2, 279--366.

\bibitem{SV-cherednik-aha}
O.~Schiffmann and E.~Vasserot,
\emph{Cherednik algebras, $W$-algebras and the equivariant cohomology of the moduli space of instantons on $\mathbb{A}^2$},
Publ.\ Math.\ IHES \textbf{118} (2013), 213--342.

\bibitem{MO}
D.~Maulik and A.~Okounkov,
\emph{Quantum groups and quantum cohomology},
Ast\'erisque \textbf{408} (2019).

\bibitem{RSYZ}
M.~Rapcak, Y.~Soibelman, Y.~Yang, and G.~Zhao,
\emph{Cohomological Hall algebras, vertex algebras, and instantons},
Comm.\ Math.\ Phys.\ \textbf{376} (2020), no.~3, 1803--1873.

\bibitem{KS-stability}
M.~Kontsevich and Y.~Soibelman,
\emph{Stability structures, motivic Donaldson--Thomas invariants and cluster transformations},
arXiv:0811.2435.

\bibitem{KS-coha}
M.~Kontsevich and Y.~Soibelman,
\emph{Cohomological Hall algebra, exponential Hodge structures and motivic Donaldson--Thomas invariants},
Commun.\ Number Theory Phys.\ \textbf{5} (2011), no.~2, 231--352.

\bibitem{Tsymbaliuk}
A.~Tsymbaliuk,
\emph{The affine Yangian of $\mathfrak{gl}_1$ revisited},
Adv.\ Math.\ \textbf{304} (2017), 583--645.

\bibitem{Prochazka-Rapcak}
T.~Proch\'azka and M.~Rapčák,
\emph{Webs of $W$-algebras},
J.\ High Energy Phys.\ \textbf{2018}, no.~11, 109.

\bibitem{Lurie-HA}
J.~Lurie,
\emph{Higher Algebra},
available at \texttt{https://www.math.ias.edu/\~{}lurie/papers/HA.pdf}.

\bibitem{Ginzburg}
V.~Ginzburg,
\emph{Calabi--Yau algebras},
arXiv:math/0612139.

\end{thebibliography}

%% ======================================================================
\appendix
\section{The Costello--Li mechanism and the $\Omega$-background}
\label{app:costello-li}
%% ======================================================================

The Costello--Li theory of twisted holography provides the
physical mechanism underlying the $E_1 \to E_2$ enhancement.
We summarize the key points relevant to the CoHA/$G(X)$ story.

\subsection{5d NCCS and $Y^+(\widehat{\gl}_1)$}

Costello--Li consider 5d non-commutative Chern--Simons (NCCS)
theory on $\CC^2 \times \CC_\hbar$, where $\CC_\hbar$ is the
holomorphic direction along the M-theory circle.  The algebra
of local operators on a line defect in this theory is the
\emph{positive half of the affine Yangian}
$Y^+(\widehat{\gl}_1)$.

The identification proceeds as follows:
\begin{enumerate}[label=(\roman*)]
  \item The fields of 5d NCCS are a $(0,1)$-form
    $A \in \Omega^{0,1}(\CC^2 \times \CC_\hbar, \gl_N)$ with
    action $S = \int \Tr(A \wedge \bar\partial A +
    \tfrac{2}{3} A \wedge A \wedge A) \wedge \Omega$, where
    $\Omega = dz_1 \wedge dz_2 \wedge d\hbar$ is the
    holomorphic volume form.

  \item The Koszul duality of this theory (in the sense of
    Costello--Gwilliam) on the line $\{0\} \times \CC_\hbar$
    produces an associative algebra.  In the $N \to \infty$ limit,
    this is $Y^+(\widehat{\gl}_1)$.

  \item The $\mathcal{W}_{1+\infty}$ algebra at self-dual level
    $\psi = 0$ is isomorphic to $Y^+(\widehat{\gl}_1)$ (by the
    Maulik--Okounkov / Schiffmann--Vasserot theorem).
\end{enumerate}

\subsection{The $\Omega$-background provides framing}

The $\Omega$-background $(\epsilon_1, \epsilon_2)$ on $\CC^2$
breaks the $E_2$ structure of the full theory down to $E_1$:

\begin{itemize}
  \item \textbf{Without $\Omega$-background}: the theory on
    $\CC^2$ is topological (holomorphic-topological), and the
    two holomorphic directions $(z_1, z_2)$ give an $E_2$-algebra
    of operators by the Dunn additivity theorem.  The
    two $E_1$ structures from the two directions combine
    into a single $E_2$ structure.

  \item \textbf{With $\Omega$-background}: the equivariant
    parameters $(\epsilon_1, \epsilon_2)$ break the symmetry
    between the two holomorphic directions.  This selects a
    preferred $E_1$-direction (say $z_1$) and demotes $z_2$ to
    a parameter.  The resulting algebra is $E_1$ (associative),
    and this is the CoHA.
\end{itemize}

\begin{remark}[The $E_1$ mechanism: $\Omega$ breaks Dunn]
\label{rem:omega-breaks-dunn}
In the language of operads, Dunn's theorem states
$E_1 \otimes E_1 \simeq E_2$.  The $\Omega$-background
gives a non-trivial equivariant structure on one of the
two $E_1$ factors, making it no longer equivalent to the
bare $E_1$ operad.  The tensor product
$E_1^{\epsilon_1} \otimes E_1^{\epsilon_2}$ is a deformation
of $E_2$ that is generically only $E_1$.  At
$\epsilon_1 + \epsilon_2 = 0$ (the ``self-dual'' or
``unrefined'' point), a residual $\ZZ/2$-symmetry survives
and the braiding is restored.  This is why
$\mathcal{W}_{1+\infty}$ at self-dual level has enhanced
symmetry.
\end{remark}

\noindent
For the quantum vertex chiral group $G(X)$: the full $E_2$
algebra is the universal structure, and the CoHA is its
$\Omega$-deformed $E_1$ shadow.  Removing the
$\Omega$-background recovers the braiding and hence the
quantum group $R$-matrix.

\end{document}
